\begin{table}
\centering
\small
\begin{tabular}{lccccccc}
\toprule
Condition & Train loss & Train acc (\%) & Val loss & Val acc (\%) & Spikes & $\Delta$ vs curve (pp) & Spikes / neuron \\
\midrule
Ours & 0.6043 & 95.38 & 0.5893 & 96.71 & 194,747 & -- & -- \\
Ours & 0.5997 & 95.62 & 0.5899 & 96.51 & 191,493 & -- & -- \\
Ours & 0.6005 & 95.55 & 0.5952 & 96.53 & 188,208 & -- & -- \\
Ours & 0.6092 & 95.16 & 0.5905 & 96.69 & 193,048 & -- & -- \\
\cmidrule{1-8}
Mean ($n=4$) & 0.6034 & 95.428 & 0.5912 & 96.610 & 191,874 & -- & 1.4790 $\pm$ 0.0157 ($n{=}11$) \\
\midrule
w/o potential & 0.6085 & 95.18 & 0.5991 & 96.27 & 145,909 & -- & -- \\
w/o potential & 0.6111 & 94.98 & 0.5892 & 96.55 & 151,177 & -- & -- \\
w/o potential & 0.6044 & 95.31 & 0.6001 & 96.29 & 149,507 & -- & -- \\
w/o potential & 0.6028 & 95.45 & 0.6002 & 96.30 & 145,649 & -- & -- \\
\cmidrule{1-8}
Mean ($n=4$) & 0.6067 & 95.230 & 0.5971 & 96.353 & 148,061 & $-0.222$ ($3.5\sigma$, $n{=}4$) & 1.4860 ($n{=}3$) \\
\midrule
w/o final-step & 0.6065 & 95.29 & 0.5892 & 96.58 & 163,260 & -- & -- \\
w/o final-step & 0.6031 & 95.48 & 0.6030 & 96.20 & 155,422 & -- & -- \\
\cmidrule{1-8}
Mean ($n=2$) & 0.6048 & 95.385 & 0.5961 & 96.390 & 159,341 & -- & 1.4478 ($n{=}2$) \\
\midrule
w/o inversion & 0.5961 & 95.70 & 0.5945 & 96.52 & 441,044 & -- & -- \\
w/o inversion & 0.5948 & 95.79 & 0.5923 & 96.61 & 432,977 & -- & -- \\
\cmidrule{1-8}
Mean ($n=2$) & 0.5955 & 95.745 & 0.5934 & 96.565 & 437,011 & -- & 1.5300 ($n{=}2$) \\
\bottomrule
\end{tabular}
\caption{Ablation, tier A: all arms at the same $\rho=4\times10^{-3}$ (spec \S 12-4 tier A). Removing a component changes the neuron budget $R$, so the auto-solved strength moves the landing point in spike count even though $\rho$ is held fixed -- that movement is itself the result (spec \S 2 A9). $\Delta$ vs curve is the accuracy offset against the fitted accuracy-spike curve (spec \S 2 A6); only the sub-threshold term ablation has this statistic computed at this sample size, so the other two arms show \texttt{--}. Spikes / neuron is drawn from a separate neuron-level extraction (spec \S 8) whose run count differs from the accuracy columns and is cited per row; arms with $n<4$ report $n$ only, with no fabricated standard deviation.}
\label{tab:t3_ablation}
\end{table}

% GATE: T3-matched -- tier B (spike-matched at ~192K) is designed but the
% 12 runs it needs have NOT been executed and are NOT approved (spec \S 12-4).
% Targets, once run, go through the same collect_paper.py + gen_tables.py path:
%   w/o potential   (abl_novmem)  rho ~= 2.6e-3  -> target ~192K spikes
%   w/o final-step  (abl_nofinal) rho ~= 2.9e-3  -> target ~192K spikes
%   fixed $\lambda$ (abl_nolr)    lambda ~= 3.2e-7 -> target ~192K spikes
%   w/o inversion   (abl_noinv)   NOT matched -- its unmatched landing point
%     of 437,011 against the ~192K the other arms land at (at the same rho=4e-3
%     as Ours) is the evidence that the inversion is a precondition for the
%     method to work at all, not a tunable knob (spec \S 2 A5).
% Do not fill this block in without a re-approval -- see spec \S 12-4.
